\documentclass[11pt,a4paper]{article}
\usepackage[margin=1in]{geometry}
\usepackage[T1]{fontenc}
\usepackage{titlesec}
\usepackage{enumitem}
\usepackage{booktabs}
\usepackage{longtable}
\usepackage{xcolor}
\usepackage{fancyhdr}
\usepackage{hyperref}
\usepackage{parskip}
\usepackage{listings}

\hypersetup{
  colorlinks=true,
  linkcolor=blue!60!black,
  urlcolor=blue!60!black,
  citecolor=blue!60!black
}

\definecolor{codebg}{RGB}{244,244,244}
\lstset{
  backgroundcolor=\color{codebg},
  basicstyle=\ttfamily\small,
  breaklines=true,
  frame=single,
  columns=fullflexible,
  showstringspaces=false
}

\pagestyle{fancy}
\fancyhf{}
\fancyhead[L]{PetSocial --- Deployment Guide}
\fancyfoot[C]{\thepage}
\renewcommand{\headrulewidth}{0.4pt}

\title{PetSocial\\\large Deployment Guide}
\author{PetSocial Development Documentation}
\date{\today}

\begin{document}
\maketitle
\thispagestyle{fancy}
\tableofcontents

\section{Overview}
PetSocial deploys as two independent pieces: the Next.js application (hosted on
Vercel) and its data layer (a Supabase project, providing PostgreSQL, Auth, and file
storage). Vercel builds and redeploys automatically on every push to the connected Git
branch, giving continuous delivery with no separate release step.

\section{Prerequisites}
\begin{itemize}[leftmargin=*]
  \item Node.js 20 or later and npm.
  \item A Supabase account and project (free tier is sufficient for development).
  \item A GitHub account and a Vercel account (Vercel can authenticate directly against
        GitHub).
\end{itemize}

\section{Supabase Project Setup}

\subsection{Database connection strings}
Supabase exposes two relevant Postgres connection modes, and the project uses both for
different purposes:

\begin{itemize}[leftmargin=*]
  \item \textbf{Transaction pooler} (port 6543) --- used by the running application at
        runtime, because it handles many short-lived connections well, as is typical in
        a serverless deployment. Requires the query parameter
        \texttt{pgbouncer=true}, since transaction-mode pooling does not support
        prepared statements.
  \item \textbf{Session pooler} (port 5432) --- used only for running migrations, since
        transaction-mode pooling does not support the schema-changing statements
        Prisma Migrate issues.
\end{itemize}

Both are found in the Supabase dashboard under \textbf{Settings $\to$ Database
$\to$ Connection string}.

\subsection{Environment variables}
\begin{longtable}{@{}p{4.6cm}p{9.4cm}@{}}
\toprule
\textbf{Variable} & \textbf{Purpose} \\
\midrule
\endhead
\texttt{DATABASE\_URL} & Transaction pooler connection string, with
  \texttt{?pgbouncer=true} appended. Used by the application at runtime. \\
\texttt{DIRECT\_URL} & Session pooler connection string. Used only by
  \texttt{prisma migrate}. \\
\texttt{NEXT\_PUBLIC\_SUPABASE\_URL} & The project's API URL, from
  \textbf{Settings $\to$ API}. \\
\texttt{NEXT\_PUBLIC\_SUPABASE\_PUBLISHABLE\_KEY} & The project's publishable
  (anonymous) API key, from the same page. \\
\bottomrule
\end{longtable}

Copy \texttt{.env.example} to \texttt{.env} for local development and fill in the four
values above; \texttt{.env} and \texttt{.env.local} are both git-ignored and must never
be committed.

\subsection{Storage bucket and policies}
Every user-uploaded image (post photos, avatars, adoption and marketplace listing
photos) is stored in a single public bucket named \texttt{post-media}.

\begin{enumerate}[leftmargin=*]
  \item In the Supabase dashboard, go to \textbf{Storage $\to$ New bucket}, name it
        \texttt{post-media}, and mark it \textbf{public}.
  \item In \textbf{SQL Editor}, run:
\end{enumerate}

\begin{lstlisting}
create policy "Allow authenticated uploads to post-media"
on storage.objects for insert
to authenticated
with check (bucket_id = 'post-media');

create policy "Allow public read of post-media"
on storage.objects for select
to public
using (bucket_id = 'post-media');
\end{lstlisting}

The dashboard's built-in policy templates were found, during development, to be scoped
incorrectly for this use case; the explicit SQL above is the known-good version.

\subsection{Authentication settings}
\begin{itemize}[leftmargin=*]
  \item For local development, email confirmation can be turned off under
        \textbf{Authentication $\to$ Sign In / Providers $\to$ Email}, so that signing
        up logs a user in immediately without needing to click a confirmation link.
  \item \textbf{This must be turned back on before any public deployment.}
  \item Supabase's default built-in email sender has a very low rate limit (a handful
        of emails per hour). Before relying on real confirmation or password-reset
        emails in production, configure a custom SMTP provider under
        \textbf{Authentication $\to$ Emails}.
\end{itemize}

\section{Local Development}
\begin{lstlisting}
npm install
npx prisma migrate dev --name init   # only needed once, against a fresh database
npm run dev
\end{lstlisting}
The application is then available at \texttt{http://localhost:3000}.

\section{Database Migrations}
Every schema change is applied through Prisma Migrate, never by hand-editing the
database:
\begin{lstlisting}
npx prisma migrate dev --name <short_description>
\end{lstlisting}
This command reads \texttt{DIRECT\_URL} (via \texttt{prisma.config.ts}), creates a new
timestamped folder under \texttt{prisma/migrations/}, applies it, and regenerates the
Prisma Client. Migration folders are committed to version control; the generated
client is not (see below).

\section{The Generated Prisma Client}
The Prisma Client is generated into \texttt{src/generated/prisma/} and is
\textbf{deliberately excluded from version control} via \texttt{.gitignore}, since it
is build output, not source. This has one important consequence: any environment that
checks out the repository fresh --- including Vercel --- must regenerate the client
before building, or the build fails with:
\begin{lstlisting}
Module not found: Can't resolve '@/generated/prisma/client'
\end{lstlisting}
This is handled by a \texttt{postinstall} script in \texttt{package.json}:
\begin{lstlisting}
"scripts": {
  "postinstall": "prisma generate"
}
\end{lstlisting}
which npm runs automatically after every \texttt{npm install}, including the one
Vercel performs at the start of every build. If this script is ever removed, every
fresh deployment will fail with the error above.

\section{GitHub}
\begin{enumerate}[leftmargin=*]
  \item Create a new, empty repository on GitHub (no README, license, or
        \texttt{.gitignore} --- the project already has its own).
  \item From the project directory:
\begin{lstlisting}
git remote add origin https://github.com/<you>/<repo>.git
git branch -M main
git push -u origin main
\end{lstlisting}
\end{enumerate}

\section{Vercel}
\begin{enumerate}[leftmargin=*]
  \item At \url{https://vercel.com/new}, sign in (GitHub sign-in is simplest) and click
        \textbf{Import} next to the repository.
  \item Vercel auto-detects the Next.js framework; the default build settings do not
        need to change.
  \item Before deploying, add the four environment variables listed in
        \S~3.2 under \textbf{Environment Variables}.
  \item Click \textbf{Deploy}. Every subsequent push to the connected branch triggers a
        new build and deployment automatically --- there is no manual redeploy step for
        ordinary changes.
\end{enumerate}

\section{Post-Deployment Checklist}
\begin{itemize}[leftmargin=*]
  \item Re-enable Supabase email confirmation and configure custom SMTP (\S~3.4).
  \item Rotate the database password if it was ever shared outside of the Supabase
        dashboard or a properly secured secrets manager.
  \item Replace the in-memory rate limiter (\texttt{src/lib/rate-limit.ts}) with a
        shared store (e.g.\ Upstash Redis) if the deployment runs across more than one
        serverless instance; see the System Architecture document, \S~``Rate
        Limiting''.
  \item Confirm the Storage bucket policies from \S~3.3 are present --- they are
        project-level Supabase configuration and are not recreated by any code in this
        repository.
\end{itemize}

\end{document}
